%%=============================================================================
%% CAPITULO 4 -- ANALISE E DISCUSSAO
%% Versao publica: apresenta apenas o marco normativo publico. Os resultados
%% da pesquisa nao sao divulgados nesta versao.
%%=============================================================================

\section{Análise e discussão}

A análise da conformidade da solução toma como referência o marco regulatório
efetivamente aplicável à PGDF: a Lei Geral de Proteção de Dados
\parencite{brasil2018lgpd}, a Portaria PGDF nº~160/2025, que aprova o Guia de
Inteligência Artificial Generativa \parencite{df2025portaria160}, e o Decreto
Distrital nº~48.901/2026, que institui a Política de Governança de
Inteligência Artificial do Distrito Federal \parencite{df2026decreto48901}. A
Tabela~\ref{tab:marco-normativo} sintetiza o papel de cada norma na pesquisa.

\begin{table}[htb]
  \caption{Marco normativo da pesquisa e papel de cada norma}
  \label{tab:marco-normativo}
  \centering
  \begin{tabular}{lll}
    \toprule
    \textbf{Norma} & \textbf{Situação} & \textbf{Papel na pesquisa} \\
    \midrule
    Lei nº 13.709/2018 (LGPD)        & vigente      & vinculante \\
    Portaria PGDF nº 160/2025        & vigente      & vinculante \\
    Decreto Distrital nº 48.901/2026 & vigente      & vinculante \\
    Resolução CNJ nº 615/2025        & vigente      & referência analógica \\
    Projeto de Lei nº 2.338/2023     & em tramitação & parâmetro prospectivo \\
    \bottomrule
  \end{tabular}
  \fonte{Elaborada pela autora (2026).}
\end{table}

A Resolução CNJ nº~615/2025, que disciplina o desenvolvimento, a utilização e
a governança de soluções de inteligência artificial no Poder Judiciário
\parencite{cnj2025res615}, é dirigida a outro Poder e não vincula a PGDF; seus
parâmetros são mobilizados nesta pesquisa apenas como referência normativa
comparativa e analógica. Cumpre distinguir, ainda, o marco legal de
inteligência artificial \parencite{brasil2023pl2338}, que se encontra em
tramitação no Congresso Nacional e, portanto, não constitui norma vigente,
devendo ser tratado como parâmetro prospectivo.
